\documentclass[12pt,a4paper]{article}

% --- PAQUETES ---
\usepackage[utf8]{inputenc}
\usepackage[T1]{fontenc}
\usepackage[spanish]{babel}
\usepackage{amsmath, amssymb, mathtools} % Paquetes matemáticos principales
\usepackage{geometry}
\usepackage{parskip} % Párrafos separados por espacio, no por sangría

% --- MÁRGENES ---
\geometry{
 a4paper,
 left=2.5cm,
 right=2.5cm,
 top=2.5cm,
 bottom=2.5cm
}

% --- COMANDOS PERSONALIZADOS ---
% Simplifican la escritura de comandos matemáticos comunes
\newcommand{\C}{\mathbb{C}}
\newcommand{\Poly}{\mathbb{P}[z]}
\newcommand{\D}{\mathcal{D}}
\newcommand{\Vdos}{V_2}

% Normas
\newcommand{\normS}[1]{\left\| #1 \right\|_S^2}
\newcommand{\normLdos}[1]{\left\| #1 \right\|_{L^2(\mu)}^2}

% Operadores y misceláneos
\newcommand{\ket}[1]{\left| #1 \right|}
\DeclareMathOperator{\supp}{supp}
\DeclareMathOperator{\codim}{codim}

% --- TÍTULO ---
\title{Cota $\sigma_2$}
\date{} % Sin fecha

% --- DOCUMENTO ---
\begin{document}

\maketitle

Demostramos que el subespacio $V_2 = \{ p \in \Poly : p(0) = 0 \}$ es un candidato válido, lo que significa que la norma del operador restringido a él es finita.

La norma al cuadrado del operador restringido es:
$$ || \D \rvert_{\Vdos} ||_S^2 = \sup_{p \in \Vdos, p \ne 0} \frac{\normS{\D(p)}}{\normS{p}} $$

Analizamos el numerador y el denominador por separado para un $p \in \Vdos$.

Para $p \in \Vdos$, $p(0)=0$. La norma de $p$ es:
$$ \normS{p} = \int_{S_R(C)} |p(z)|^2 d\mu(z) + |p'(0)|^2 $$
Dado que $p'(0)$ no es necesariamente cero, ambos términos de la norma de Sobolev se mantienen.

La imagen de $p$ bajo $\D$ es $\D(p) = zp$. Su norma de Sobolev es:
$$ \normS{\D(p)} = \int_{S_R(C)} |z p(z)|^2 d\mu(z) + \ket{(zp)' \rvert_{z=0}}^2 $$
El término atómico es el punto crítico. Usando la regla del producto para la derivada:
$$ (zp)' \rvert_{z=0} = [p(z) + z p'(z)]_{z=0} = p(0) + 0 \cdot p'(0) = p(0) $$
Dado que $p \in \Vdos$, se cumple por definición que $p(0) = 0$. Por lo tanto, el término atómico de la imagen se anula:
$$ \ket{(zp)' \rvert_{z=0}}^2 = |p(0)|^2 = 0 $$
El numerador se simplifica a la parte integral:
$$ \normS{\D(p)} = \int_{S_R(C)} |z p(z)|^2 d\mu(z) $$

Sustituimos los resultados del numerador y denominador en la fórmula del supremo:
$$ || \D \rvert_{\Vdos} ||_S^2 = \sup_{p \in \Vdos} \frac{\int_{S_R(C)} |z p(z)|^2 d\mu(z)}{\int_{S_R(C)} |p(z)|^2 d\mu(z) + |p'(0)|^2} $$
Para encontrar una cota superior, notamos que el término $|p'(0)|^2$ en el denominador es no negativo ($> 0$). Por lo tanto, el denominador es siempre mayor o igual que la parte integral:
$$ \int_{S_R(C)} |p|^2 d\mu + |p'(0)|^2 \ge \int_{S_R(C)} |p|^2 d\mu = \normLdos{p} $$
Al usar un denominador más pequeño (o igual) en la fracción, el valor de la fracción se hace más grande (o igual), lo que nos da una cota superior válida:
$$ || \D \rvert_{\Vdos} ||_S^2 \le \sup_{p \in \Vdos} \frac{\int_{S_R(C)} |z p|^2 d\mu}{\int_{S_R(C)} |p|^2 d\mu} = \sup_{p \in \Vdos} \frac{\normLdos{zp}}{\normLdos{p}} $$
El término de la derecha es la norma al cuadrado del operador de multiplicación por $z$ en el espacio $L^2(\mu)$, restringido a $\Vdos$. Esta norma es, como máximo, la norma del operador en todo el espacio $L^2(\mu)$:
$$ || \D \rvert_{\Vdos} ||_S^2 \le ||\D||_{L^2(\mu)}^2 = \left( \sup_{z \in \supp(\mu)} |z| \right)^2 $$
El soporte es $S_R(C)$. El punto en esta circunferencia más alejado del origen es $|C|+R$ (por la desigualdad triangular).
$$ ||\D||_{L^2(\mu)}^2 = (|C| + R)^2 $$

Hemos demostrado que la norma del operador restringido a $\Vdos = \ker(p(0))$ es finita y está acotada:
$$ || \D \rvert_{\Vdos} ||_S^2 \le (|C| + R)^2 < \infty $$
Esto confirma que $\Vdos$ es un candidato válido que proporciona una cota superior finita para $\lim \sigma_2$.

\end{document}